\documentclass[a4paper,11pt]{article}
\def\email#1{{\tt#1}}

\def\N{\mbox{\rm{I}\hspace{-.15em}\rm{N}}}
\def\Z{\mbox{\rm{Z}\hspace{-.28em}\rm{Z}}}
\def\RR{\mbox{\rm{I}\hspace{-.15em}\rm{R}}}
\def\C{\mbox{\rm{l}\hspace{-.50em}\rm{C}}}
\def\Q{\mbox{\rm{l}\hspace{-.50em}\rm{Q}}}
\def\F{\mbox{\rm{I}\hspace{-.20em}\rm{F}}}

\usepackage{amssymb,amsmath}
\usepackage{url}
\usepackage{fullpage}
\begin{document}

\newtheorem{theo}{Theorem}
\newtheorem{coro}{Corollary}[theo]
\newtheorem{lemme}{Lemma}
\renewcommand{\thecoro}{\thetheo.\arabic{coro}}



\parindent=0cm


\section*{\center IN100 Initiation au Python\\
Contrôle continu 2}

\vspace{1 cm}

\begin{table}[h!]
\centering
\renewcommand{\arraystretch}{1.4} % augmente la hauteur des lignes
\begin{tabular}{|p{10cm}|p{3cm}|p{2cm}|}
\hline
\textbf{Étudiant :} & \textbf{Heure Début }  & \textbf{Note (/5) }  \\
& & \\
\hline
\textbf{Correcteur :} & \textbf{Heure Fin } & \\
& & \\
\hline

\hline
\end{tabular}
\end{table}


\subsection*{Consignes:} Vous avez 20 minutes pour coder la solution. Seules les types et les instructions vues en cours sont autorisées. Tout usage d'une fonction 
venant d'une bibliothèque qui n'a pas été vue en cours ne sera pris en compte. 


\subsection*{Sujet}

Un billet de loterie est donné par 5 nombres entiers compris entre 1 et 54. 


\begin{itemize}
\item[1.]  En utilisant les notions vues en cours, donnez une fonction permettant de générer un billet quelconque pour un joueur. 
Nous rappelons que la fonction \texttt{random.randint(a,b)} permet de générer des nombres aléatoires entre \texttt{a} et \texttt{b}. 

Exemple d'usage: 
\begin{verbatim}
import random

print(random.randint(3,9))

\end{verbatim}

\item[2.] Choisissez un billet gagnant pour votre loterie. \'Ecrire une fonction permettant de vérifier, pour un billet quelconque, si celui-ci est gagnant. 
\end{itemize}



\begin{table}[h!]
\centering
\renewcommand{\arraystretch}{1.4} % augmente la hauteur des lignes
\begin{tabular}{|p{15cm}|}
\hline
\textbf{Remarques étudiant :}    \\
 \\
 \\
\hline
\textbf{Remarques correcteur :}  \\
 \\
 \\
\hline
\end{tabular}
\end{table}

\newpage


\section*{\center IN100 Initiation au Python\\
Contrôle continu 2}


\begin{table}[h!]
\centering
\renewcommand{\arraystretch}{1.4} % augmente la hauteur des lignes
\begin{tabular}{|p{10cm}|p{3cm}|p{2cm}|}
\hline
\textbf{Étudiant :} & \textbf{Heure Début }  & \textbf{Note (/5) }  \\
& & \\
\hline
\textbf{Correcteur :} & \textbf{Heure Fin } & \\
& & \\
\hline

\hline
\end{tabular}
\end{table}

\subsection*{Consignes:} Vous avez 20 minutes pour coder la solution. Seules les types et les instructions vues en cours sont autorisées. Tout usage d'une fonction 
venant d'une bibliothèque qui n'a pas été vue en cours ne sera pris en compte. 


\subsection*{Sujet:}


\begin{itemize}
\item[1.] \'Ecrire une fonction qui calcule le mot le plus long dans une chaîne de caractères et le renvoie. 
Par exemple, sur la chaîne suivante 
\begin{verbatim}
chaine="Paris e beau"
\end{verbatim}
la fonction renverra \texttt{"Paris"}.  On suppose que les chaînes contiennent que les lettres de l'alphabet et le caractère \texttt{" "}. 

\item[2.] On  donne une liste contenant que des chaînes de caractères. 
\'Ecrire une fonction qui à partir de la liste, calcule le mot le plus long dans ces chaînes et renvoie l'indice dans la liste de la chaîne  contenant ce mot ainsi que le mot. 
\end{itemize}

Par exemple sur la liste:
\begin{verbatim}

capitales=["Paris e beau","A Madrid il fait chaud", "Londres est pluvieux", 
"Bruxelles est magnifique", "A Prague il fait froid"]

\end{verbatim}
votre fonction devra renvoyer \texttt{(3,magnifique)}. 

\begin{table}[h!]
\centering
\renewcommand{\arraystretch}{1.4} % augmente la hauteur des lignes
\begin{tabular}{|p{15cm}|}
\hline
\textbf{Remarques étudiant :}    \\
 \\
 \\
\hline
\textbf{Remarques correcteur :}  \\
 \\
 \\
\hline
\end{tabular}
\end{table}


\newpage


\section*{\center IN100 Initiation au Python\\
Contrôle continu 2}


\begin{table}[h!]
\centering
\renewcommand{\arraystretch}{1.4} % augmente la hauteur des lignes
\begin{tabular}{|p{10cm}|p{3cm}|p{2cm}|}
\hline
\textbf{Étudiant :} & \textbf{Heure Début }  & \textbf{Note (/5) }  \\
& & \\
\hline
\textbf{Correcteur :} & \textbf{Heure Fin } & \\
& & \\
\hline

\hline
\end{tabular}
\end{table}

\subsection*{Consignes:} Vous avez 20 minutes pour coder la solution. Seules les types et les instructions vues en cours sont autorisées. Tout usage d'une fonction 
venant d'une bibliothèque qui n'a pas été vue en cours ne sera pris en compte. 

\subsection*{Sujet:}

Un horaire de train La Défense - Montreuil est donné sous forme d'un tuple indiquant l'heure de départ et l'heure d'arrivée. Exemple \texttt{(14h34, 14h47)}. 

\begin{itemize}
\item[1.] \'Ecrire une fonction qui prend en entrée un tuple indiquant les horaires d'un train et retourne la durée du trajet en minutes. 
\item[2.] Bidule doit arriver à temps à son cours qui commence à une heure donnée. \'Ecrire une fonction \texttt{calcul\_depart} qui prend en argument la liste des horaires 
des trains  de la journée et l'heure de début du cours et qui renvoie la liste des trains possibles pour que Bidule  arrive à l'heure en cours. 

Vous pouvez tester votre fonction sur la liste suivante: 
\begin{verbatim}
horaires=[(9h50,10h12),(11h36,12h00), (11h47,12h13),(13h47, 14h12)]
calcul_depart(horaires, '13h50')
\end{verbatim}
\end{itemize}


\begin{table}[h!]
\centering
\renewcommand{\arraystretch}{1.4} % augmente la hauteur des lignes
\begin{tabular}{|p{15cm}|}
\hline
\textbf{Remarques étudiant :}    \\
 \\
 \\
\hline
\textbf{Remarques correcteur :}  \\
 \\
 \\
\hline
\end{tabular}
\end{table}


\newpage

\section*{\center IN100 Initiation au Python\\
Contrôle continu 2}


\begin{table}[h!]
\centering
\renewcommand{\arraystretch}{1.4} % augmente la hauteur des lignes
\begin{tabular}{|p{10cm}|p{3cm}|p{2cm}|}
\hline
\textbf{Étudiant :} & \textbf{Heure Début }  & \textbf{Note (/5) }  \\
& & \\
\hline
\textbf{Correcteur :} & \textbf{Heure Fin } & \\
& & \\
\hline

\hline
\end{tabular}
\end{table}

\subsection*{Consignes:} Vous avez 20 minutes pour coder la solution. Seules les types et les instructions vues en cours sont autorisées. Tout usage d'une fonction 
venant d'une bibliothèque qui n'a pas été vue en cours ne sera pris en compte. 

\subsection*{Sujet:}

On souhaite écrire un programme qui analyse une liste de prénoms d’élèves d’une classe.

\begin{itemize}
\item[1.] Écrire une fonction \texttt{saisir\_prenoms} qui :

\begin{itemize}
\item Demande à l’utilisateur de saisir des prénoms un par un,
\item S’arrête lorsque l’utilisateur saisit une chaîne vide (""),
\item Retourne la liste des prénoms saisis.
\end{itemize}
\item[2.]  \'Ecrire une fonction \texttt{prenom\_le\_plus\_frequent} qui prend en argument la liste de prenoms
et retourne le prénom qui apparaît le plus souvent dans la liste.
\end{itemize}
\vspace{1 cm}

\begin{table}[h!]
\centering
\renewcommand{\arraystretch}{1.4} % augmente la hauteur des lignes
\begin{tabular}{|p{15cm}|}
\hline
\textbf{Remarques étudiant :}    \\
 \\
 \\
\hline
\textbf{Remarques correcteur :}  \\
 \\
 \\
\hline
\end{tabular}
\end{table}

\newpage

\section*{\center IN100 Initiation au Python\\
Contrôle continu 2}


\begin{table}[h!]
\centering
\renewcommand{\arraystretch}{1.4} % augmente la hauteur des lignes
\begin{tabular}{|p{10cm}|p{3cm}|p{2cm}|}
\hline
\textbf{Étudiant :} & \textbf{Heure Début }  & \textbf{Note (/5) }  \\
& & \\
\hline
\textbf{Correcteur :} & \textbf{Heure Fin } & \\
& & \\
\hline

\hline
\end{tabular}
\end{table}

\subsection*{Consignes:} Vous avez 20 minutes pour coder la solution. Seules les types et les instructions vues en cours sont autorisées. Tout usage d'une fonction 
venant d'une bibliothèque qui n'a pas été vue en cours ne sera pris en compte. 


\subsection*{Sujet:}

On souhaite écrire un programme qui aide à analyser les ventes journalières d’un petit magasin.

\begin{itemize}
\item[1.] Écrire une fonction \texttt{saisir\_ventes} qui :
\begin{itemize}
\item Demande à l’utilisateur de saisir les montants des ventes (en euros) d’une journée.
\item Continue la saisie jusqu’à ce que l’utilisateur entre 0.
\item Retourne la liste des montants saisis.
\end{itemize}
\item[2.]  \'Ecrire une fonction \texttt{total\_ventes} qui  prend en entrée la liste des montants et retourne la somme de toutes les ventes. 
\item[3.] \'Ecrire une fonction \texttt{superieures\_a\_moyenne} qui prend en entrée la liste de montants et  retourne la liste des ventes supérieures à la moyenne. 
\end{itemize}

\vspace*{0.2 cm}
\begin{table}[h!]
\centering
\renewcommand{\arraystretch}{1.4} % augmente la hauteur des lignes
\begin{tabular}{|p{15cm}|}
\hline
\textbf{Remarques étudiant :}    \\
 \\
 \\
\hline
\textbf{Remarques correcteur :}  \\
 \\
 \\
\hline
\end{tabular}
\end{table}





\end{document}
